\documentclass[aspectratio=169]{beamer}
\usetheme{Madrid}
\usecolortheme{default}
\usepackage[utf8]{inputenc}
\usepackage{amsmath}
\usepackage{amsfonts}
\usepackage{amssymb}
\usepackage{enumerate}

\setbeamertemplate{navigation symbols}{}
\setbeamerfont{block body}{size=\small}
\setbeamerfont{alertblock body}{size=\small}

\title{CRISC Mixed Practice Exam III -- 100 Real-Life Questions}
\subtitle{Weighted to the Four CRISC Domains -- Practical Scenarios}
\author{CRISC Mastery Series}
\date{\today}

% ------------------------------------------------------------
% \hq{number}{topic}{question}{options}{letter}{explanation}{takeaway}
% Creates one question frame and one explanation frame.
% ------------------------------------------------------------
\newcommand{\hq}[7]{%
	\begin{frame}{Q#1: #2}
		\begin{block}{Question}
			#3
		\end{block}
		\begin{itemize}
			#4
		\end{itemize}
	\end{frame}
	\begin{frame}{Q#1: Explanation}
		\begin{block}{Correct Answer: #5}
			\begin{itemize}
				#6
			\end{itemize}
		\end{block}
		\begin{alertblock}{Key Takeaway}
			#7
		\end{alertblock}
	\end{frame}
}

\begin{document}

	\begin{frame}
		\titlepage
	\end{frame}

	\begin{frame}{Instructions}
		\begin{alertblock}{How to Use This Deck}
			Think like a risk practitioner in a real organization.\\
			Focus on business context, ownership, and practical decision-making.\\
			The answer and a key takeaway are on the slide following each question.
		\end{alertblock}
	\end{frame}

	\begin{frame}{Outline}
		\tableofcontents
	\end{frame}

	% ============================================================
	\section{Governance (26 percent of the exam)}
	% ============================================================

	\hq{1}{Risk Decision Log}
	{Leaders make verbal risk decisions in meetings with no record, and later disagree about what was agreed. What is the \textbf{BEST} fix?}
	{\item[\textbf{(A)}] Rely on memory
		\item[\textbf{(B)}] Maintain a decision log capturing the risk, options, decision, owner, and date
		\item[\textbf{(C)}] Stop holding risk meetings
		\item[\textbf{(D)}] Assign the practitioner to remember}
	{B}
	{\item Records support accountability.
		\item Decisions can be reviewed when conditions change.
		\item Owners and dates prevent confusion.}
	{If it is not recorded, it was not decided.}

	\hq{2}{Board Risk Committee Composition}
	{Which composition \textbf{BEST} supports an effective board risk committee?}
	{\item[\textbf{(A)}] Only the management team
		\item[\textbf{(B)}] Members with relevant risk and business expertise and sufficient independence from management
		\item[\textbf{(C)}] Only technical staff
		\item[\textbf{(D)}] Members chosen for availability alone}
	{B}
	{\item Independence supports objective oversight.
		\item Expertise enables informed challenge.
		\item Mix of skills covers strategic and technical risk.}
	{Oversight needs independence and expertise.}

	\hq{3}{AI Governance Committee}
	{The organization wants to govern the use of AI. Which structure is \textbf{MOST} appropriate?}
	{\item[\textbf{(A)}] One developer deciding alone
		\item[\textbf{(B)}] No oversight to encourage speed
		\item[\textbf{(C)}] Only the marketing team
		\item[\textbf{(D)}] A cross-functional committee with legal, risk, security, privacy, and business representation and clear decision rights}
	{D}
	{\item AI risks span many disciplines.
		\item Decision rights prevent gaps.
		\item Business participation keeps it practical.}
	{AI governance is a team sport.}

	\hq{4}{Model Risk}
	{A statistical model drives pricing decisions. Which governance activity \textbf{BEST} manages model risk?}
	{\item[\textbf{(A)}] Trust the developer's description
		\item[\textbf{(B)}] Validate once and never again
		\item[\textbf{(C)}] Independent validation, documented assumptions, ongoing performance monitoring, and defined ownership
		\item[\textbf{(D)}] Hide model limitations}
	{C}
	{\item Models can be wrong or drift.
		\item Independent validation provides assurance.
		\item Monitoring detects degradation.}
	{Validate models independently and keep watching them.}

	\hq{5}{Climate Risk Governance}
	{Regulators and investors ask about climate-related risks to operations. What should governance do?}
	{\item[\textbf{(A)}] Treat it as marketing only
		\item[\textbf{(B)}] Ignore the question
		\item[\textbf{(C)}] Leave it to facilities alone
		\item[\textbf{(D)}] Assign ownership, integrate climate risks into the enterprise risk process, and report them to the board}
	{D}
	{\item Climate can affect operations and supply chains.
		\item Integration ensures consistent assessment.
		\item Board visibility supports oversight.}
	{Emerging risks belong in the enterprise process.}

	\hq{6}{Quantified Cyber Risk Reporting}
	{The board asks for cyber risk in terms it can compare with other business risks. What is the \textbf{BEST} approach?}
	{\item[\textbf{(A)}] Provide only technical vulnerability counts
		\item[\textbf{(B)}] Provide a list of tools
		\item[\textbf{(C)}] Translate cyber scenarios into estimated financial impact ranges and compare them to appetite
		\item[\textbf{(D)}] Avoid quantification entirely}
	{C}
	{\item Financial terms allow comparison.
		\item Ranges express uncertainty.
		\item Appetite gives context for decisions.}
	{Speak to the board in business terms.}

	\hq{7}{Culture After a Breach}
	{After a breach caused by human error, management wants to dismiss the employee. What is the \textbf{BEST} governance perspective?}
	{\item[\textbf{(A)}] Punish without investigation
		\item[\textbf{(B)}] Hide the cause
		\item[\textbf{(C)}] Focus on root causes and process or control weaknesses, applying accountability proportionately and fairly
		\item[\textbf{(D)}] Blame the vendor automatically}
	{C}
	{\item Root causes usually include system factors.
		\item Fear reduces reporting.
		\item Fair accountability supports learning.}
	{Learn from the failure; do not just blame.}

	\hq{8}{Security Champions}
	{What is the benefit of a security champions program in business units?}
	{\item[\textbf{(A)}] It embeds security awareness and local ownership, improving communication between security and the business
		\item[\textbf{(B)}] It replaces the security team
		\item[\textbf{(C)}] It removes policies
		\item[\textbf{(D)}] It lowers training needs to zero}
	{A}
	{\item Champions are trusted local voices.
		\item They help translate policy into practice.
		\item They provide feedback to security.}
	{Local champions scale security culture.}

	\hq{9}{IT Portfolio Prioritization}
	{Many IT projects compete for resources. Which approach \textbf{BEST} supports governance?}
	{\item[\textbf{(A)}] First come, first served
		\item[\textbf{(B)}] Loudest sponsor first
		\item[\textbf{(C)}] Alphabetical order
		\item[\textbf{(D)}] Prioritize by strategic value, risk, cost, and resource availability using agreed criteria}
	{D}
	{\item Agreed criteria make decisions transparent.
		\item Risk and value should be balanced.
		\item Capacity is a real constraint.}
	{Prioritize on value, risk, and capacity.}

	\hq{10}{Stage Gates}
	{A project moves from design to build without reviewing risks or security requirements. Which governance mechanism would have helped?}
	{\item[\textbf{(A)}] Longer lunch breaks
		\item[\textbf{(B)}] Stage gates requiring risk and security approval before proceeding
		\item[\textbf{(C)}] Faster coding
		\item[\textbf{(D)}] More meetings without criteria}
	{B}
	{\item Gates enforce review at key points.
		\item Early detection lowers cost.
		\item Approval creates accountability.}
	{Gates stop problems from rolling forward.}

	\hq{11}{Architecture Review Board}
	{What is the main purpose of an architecture review board?}
	{\item[\textbf{(A)}] To ensure that designs follow standards and consider security, risk, and integration impacts
		\item[\textbf{(B)}] To select office furniture
		\item[\textbf{(C)}] To approve vacation requests
		\item[\textbf{(D)}] To write marketing text}
	{A}
	{\item Consistency reduces complexity.
		\item Security is reviewed early.
		\item Exceptions are decided consciously.}
	{Review architecture before it becomes legacy.}

	\hq{12}{Policy Attestation}
	{Why require employees to attest periodically that they have read and will follow key policies?}
	{\item[\textbf{(A)}] It removes the need for training
		\item[\textbf{(B)}] It reinforces awareness, creates accountability, and provides evidence for compliance
		\item[\textbf{(C)}] It guarantees compliance
		\item[\textbf{(D)}] It lowers all risk to zero}
	{B}
	{\item Attestation documents acknowledgement.
		\item It supports enforcement.
		\item It should complement training.}
	{Attestation supports awareness and accountability.}

	\hq{13}{Data Protection Officer Independence}
	{A designated data protection officer is pressured by management to approve a risky data use. What governance principle applies?}
	{\item[\textbf{(A)}] The role must obey every instruction
		\item[\textbf{(B)}] The role must be able to act independently and report to the highest management level without conflict
		\item[\textbf{(C)}] The role can be removed for disagreeing
		\item[\textbf{(D)}] The role should hide concerns}
	{B}
	{\item Independence is important for credibility.
		\item Direct access to senior management supports escalation.
		\item Conflicts of interest should be avoided.}
	{Oversight roles need independence.}

	\hq{14}{Shadow AI}
	{Employees use unapproved AI tools to speed up their work. What is the \textbf{BEST} governance response?}
	{\item[\textbf{(A)}] Provide approved tools and clear usage rules, and monitor for risky use
		\item[\textbf{(B)}] Ignore it
		\item[\textbf{(C)}] Ban the internet
		\item[\textbf{(D)}] Punish everyone}
	{A}
	{\item Bans without alternatives push use underground.
		\item Rules clarify what data is allowed.
		\item Monitoring manages residual risk.}
	{Enable the safe option and set boundaries.}

	\hq{15}{Risk Clauses in Contracts}
	{Which contract clauses \textbf{BEST} support managing third-party risk?}
	{\item[\textbf{(A)}] Only pricing
		\item[\textbf{(B)}] Only logos
		\item[\textbf{(C)}] Only delivery dates
		\item[\textbf{(D)}] Security requirements, breach notification, audit rights, subcontractor controls, and termination assistance}
	{D}
	{\item Clauses create enforceable expectations.
		\item Notification timelines matter.
		\item Exit assistance reduces lock-in.}
	{Put your risk expectations into the contract.}

	\hq{16}{Benefits Realization}
	{A project was delivered on time but has not produced the promised business benefits. Which governance practice was weak?}
	{\item[\textbf{(A)}] Office decoration
		\item[\textbf{(B)}] Code formatting
		\item[\textbf{(C)}] Benefits realization management, including ownership and tracking of expected value
		\item[\textbf{(D)}] Server cooling}
	{C}
	{\item Delivery is not the same as value.
		\item Owners must be accountable for benefits.
		\item Tracking enables corrective action.}
	{Governance tracks value, not just delivery.}

	\hq{17}{Resource Optimization}
	{Risk and security teams are overloaded with low-value tasks. What is the \textbf{BEST} governance response?}
	{\item[\textbf{(A)}] Add more tasks
		\item[\textbf{(B)}] Prioritize work by risk, automate routine tasks, and align capacity with objectives
		\item[\textbf{(C)}] Remove oversight
		\item[\textbf{(D)}] Ignore the workload}
	{B}
	{\item Resources should focus on high risk.
		\item Automation frees capacity.
		\item Alignment avoids waste.}
	{Spend scarce expertise where risk is highest.}

	\hq{18}{Subsidiary Governance}
	{A foreign subsidiary follows weaker security practices than the group standard. What should be done?}
	{\item[\textbf{(A)}] Assess the gap, set expectations, assign accountability, and monitor remediation within group governance
		\item[\textbf{(B)}] Ignore it because it is overseas
		\item[\textbf{(C)}] Close the subsidiary immediately without analysis
		\item[\textbf{(D)}] Hide the issue from the board}
	{A}
	{\item Group risk includes subsidiaries.
		\item Local laws may require adaptation.
		\item Monitoring ensures progress.}
	{The group is exposed to the subsidiary's risk.}

	\hq{19}{Group vs Local Policy}
	{Local law conflicts with a global policy requirement. What is the \textbf{BEST} approach?}
	{\item[\textbf{(A)}] Ignore local law
		\item[\textbf{(B)}] Ignore the global policy
		\item[\textbf{(C)}] Let each employee decide
		\item[\textbf{(D)}] Seek legal advice, document the conflict, and adopt a compliant local variant approved by governance}
	{D}
	{\item Legal compliance must come first.
		\item Documentation explains deviations.
		\item Governance approval maintains consistency.}
	{Resolve conflicts through legal and governance.}

	\hq{20}{Data Sovereignty}
	{New regulations require certain data to remain within national borders. What should governance do?}
	{\item[\textbf{(A)}] Ignore until audited
		\item[\textbf{(B)}] Identify affected data, assess architecture and vendors, and plan changes to comply
		\item[\textbf{(C)}] Move all data abroad
		\item[\textbf{(D)}] Delete the data}
	{B}
	{\item Data location is a compliance requirement.
		\item Vendors may also be affected.
		\item Planning avoids rushed changes.}
	{Know where regulated data lives.}

	\hq{21}{Legal Hold}
	{A lawsuit is anticipated. How should the organization treat relevant records?}
	{\item[\textbf{(A)}] Delete them quickly
		\item[\textbf{(B)}] Preserve them under a legal hold, suspending routine deletion for the affected data
		\item[\textbf{(C)}] Continue normal deletion
		\item[\textbf{(D)}] Rewrite them}
	{B}
	{\item Destroying evidence creates serious legal risk.
		\item Legal counsel directs the scope.
		\item Hold procedures must be communicated.}
	{Preserve first, ask questions later.}

	\hq{22}{Stakeholder Mapping}
	{Why map stakeholders for risk communication?}
	{\item[\textbf{(A)}] To avoid communication
		\item[\textbf{(B)}] To tailor communication to their interests, influence, and information needs
		\item[\textbf{(C)}] To create conflicts
		\item[\textbf{(D)}] To hide information}
	{B}
	{\item Different audiences need different messages.
		\item Influence affects decisions.
		\item Early engagement builds support.}
	{Know who needs to know what.}

	\hq{23}{Facilitating Risk Workshops}
	{A risk workshop is dominated by one senior person. What is the \textbf{BEST} facilitation technique?}
	{\item[\textbf{(A)}] Let the senior person decide everything
		\item[\textbf{(B)}] End the workshop
		\item[\textbf{(C)}] Use structured input methods such as anonymous voting or round-robin contributions
		\item[\textbf{(D)}] Ignore the imbalance}
	{C}
	{\item Dominance biases results.
		\item Structured methods encourage all views.
		\item Better input improves quality.}
	{Structure the discussion to hear everyone.}

	\hq{24}{Consolidating Registers}
	{Each department keeps its own risk register in a different format. What is the \textbf{BEST} improvement?}
	{\item[\textbf{(A)}] Keep formats different
		\item[\textbf{(B)}] Adopt a common structure and criteria so that risks can be consolidated and compared enterprise-wide
		\item[\textbf{(C)}] Remove registers
		\item[\textbf{(D)}] Allow only IT to keep one}
	{B}
	{\item Common structure enables aggregation.
		\item Criteria ensure comparability.
		\item Enterprise view supports decisions.}
	{Standardize to consolidate.}

	\hq{25}{AI Incident Accountability}
	{An AI system causes a harmful decision. Who remains accountable?}
	{\item[\textbf{(A)}] The algorithm
		\item[\textbf{(B)}] The vendor only
		\item[\textbf{(C)}] The organization and the designated business owner of the process, not the algorithm or vendor alone
		\item[\textbf{(D)}] No one}
	{C}
	{\item Accountability stays with people.
		\item Owners must oversee automated decisions.
		\item Contracts may share costs but not accountability.}
	{Algorithms are not accountable; people are.}

	\hq{26}{Reporting Concerns Properly}
	{A practitioner believes a material risk is being concealed from the board. What should they do?}
	{\item[\textbf{(A)}] Ignore it
		\item[\textbf{(B)}] Post it publicly at once
		\item[\textbf{(C)}] Delete the evidence
		\item[\textbf{(D)}] Escalate through designated channels such as the audit committee, document concerns, and follow the professional code}
	{D}
	{\item Escalation uses proper channels.
		\item Documentation protects integrity.
		\item Public disclosure may be a last resort under law.}
	{Escalate through proper channels.}

	% ============================================================
	\section{Risk Assessment (20 percent of the exam)}
	% ============================================================

	\hq{27}{Ransomware Scenario}
	{When building a ransomware scenario, which elements \textbf{BEST} support a credible assessment?}
	{\item[\textbf{(A)}] The CEO's favorite color
		\item[\textbf{(B)}] Entry points, backup resilience, recovery times, business impact, and detection capability
		\item[\textbf{(C)}] Only the ransom amount
		\item[\textbf{(D)}] Only the number of servers}
	{B}
	{\item Scenarios link cause and impact.
		\item Recovery capability determines impact.
		\item Detection affects dwell time.}
	{A good scenario covers entry, impact, and recovery.}

	\hq{28}{Business Email Compromise}
	{Which control \textbf{BEST} reduces losses from business email compromise payment fraud?}
	{\item[\textbf{(A)}] Out-of-band verification of payment changes and dual approval for large transfers
		\item[\textbf{(B)}] Larger email attachments
		\item[\textbf{(C)}] Colorful email templates
		\item[\textbf{(D)}] Faster payment processing without checks}
	{A}
	{\item Fraud relies on trust in email.
		\item Independent verification breaks the attack.
		\item Dual approval limits single-person risk.}
	{Verify payment changes through another channel.}

	\hq{29}{Deepfake Fraud}
	{An attacker uses a synthetic voice to impersonate an executive and requests an urgent transfer. Which response is \textbf{BEST}?}
	{\item[\textbf{(A)}] Trust the voice
		\item[\textbf{(B)}] Skip verification for executives
		\item[\textbf{(C)}] Ignore the threat
		\item[\textbf{(D)}] Require call-back verification using known numbers and enforce approval procedures regardless of urgency}
	{D}
	{\item Voice alone is no longer reliable proof.
		\item Procedures should not be bypassed.
		\item Awareness reduces success.}
	{Do not trust urgency or voice alone.}

	\hq{30}{AI Model Drift}
	{A deployed model's accuracy slowly declines as customer behavior changes. What does the assessment need?}
	{\item[\textbf{(A)}] Nothing; models never change
		\item[\textbf{(B)}] Ongoing monitoring of performance with thresholds and a retraining or review process
		\item[\textbf{(C)}] One-time testing
		\item[\textbf{(D)}] Ignoring complaints}
	{B}
	{\item Data and environments change.
		\item Thresholds trigger action.
		\item Ownership ensures follow-up.}
	{Models need monitoring after deployment.}

	\hq{31}{Data Poisoning}
	{Which risk arises when attackers manipulate the data used to train a machine learning model?}
	{\item[\textbf{(A)}] Data poisoning, which can cause the model to behave incorrectly or in a biased way
		\item[\textbf{(B)}] Disk overheating
		\item[\textbf{(C)}] Screen glare
		\item[\textbf{(D)}] Printer queues}
	{A}
	{\item Training data quality determines behavior.
		\item Controls include source validation.
		\item Monitoring may detect anomalies.}
	{Protect the data that trains the model.}

	\hq{32}{Pandemic Scenario}
	{A prolonged pandemic forces most staff to stay away from offices. Which assessment element is \textbf{MOST} relevant?}
	{\item[\textbf{(A)}] The color of the carpets
		\item[\textbf{(B)}] Printer supplies
		\item[\textbf{(C)}] Dependency on critical roles and sites, remote capability, and continuity of essential services
		\item[\textbf{(D)}] Office plants}
	{C}
	{\item Workforce availability affects operations.
		\item Remote capability reduces impact.
		\item Critical roles need backup.}
	{Assess dependence on people and places.}

	\hq{33}{Data Center Loss}
	{Assessing the loss of a primary data center, which factor is \textbf{MOST} important?}
	{\item[\textbf{(A)}] The building color
		\item[\textbf{(B)}] The cafeteria menu
		\item[\textbf{(C)}] Dependencies, recovery capability at an alternate site, and the time to restore critical services
		\item[\textbf{(D)}] Parking spaces}
	{C}
	{\item Impact depends on recoverability.
		\item Dependencies determine order.
		\item Time to restore compares to RTO.}
	{Compare recovery capability with business need.}

	\hq{34}{Cloud Provider Outage}
	{A critical service runs in a single cloud region. What should the assessment consider?}
	{\item[\textbf{(A)}] Nothing; clouds never fail
		\item[\textbf{(B)}] Only cost
		\item[\textbf{(C)}] Regional failure impact, multi-region or alternative architectures, and recovery objectives
		\item[\textbf{(D)}] Only logos}
	{C}
	{\item Cloud services can fail.
		\item Architecture determines resilience.
		\item Objectives drive design choices.}
	{Single region means single point of failure.}

	\hq{35}{Vendor Financial Health}
	{Why consider the financial stability of a critical vendor in risk assessment?}
	{\item[\textbf{(A)}] A vendor failure could disrupt services, support, and data availability
		\item[\textbf{(B)}] It does not matter
		\item[\textbf{(C)}] Only the CFO cares
		\item[\textbf{(D)}] It guarantees security}
	{A}
	{\item Insolvency can end services abruptly.
		\item Exit plans reduce impact.
		\item Escrow arrangements may help.}
	{Vendor failure is your business risk.}

	\hq{36}{Credential Stuffing}
	{Attackers use leaked passwords from other sites to log in to customer accounts. Which control \textbf{BEST} reduces this risk?}
	{\item[\textbf{(A)}] Multi-factor authentication, rate limiting, and detection of unusual login patterns
		\item[\textbf{(B)}] Shorter passwords
		\item[\textbf{(C)}] Allow unlimited attempts
		\item[\textbf{(D)}] Disable logging}
	{A}
	{\item Password reuse is common.
		\item MFA blocks reuse alone.
		\item Rate limits and detection limit automation.}
	{Assume users reuse passwords.}

	\hq{37}{SIM Swap}
	{Customers are losing accounts because attackers transfer their phone numbers. Which authentication approach reduces this risk?}
	{\item[\textbf{(A)}] Rely only on SMS codes
		\item[\textbf{(B)}] Avoid relying only on SMS codes; use app-based or hardware authenticators
		\item[\textbf{(C)}] Remove authentication
		\item[\textbf{(D)}] Use shorter codes}
	{B}
	{\item Phone numbers can be hijacked.
		\item Stronger factors are less vulnerable.
		\item Risk-based step-up can help.}
	{SMS alone is a weak second factor.}

	\hq{38}{Physical Climate Risk}
	{A data center sits in a region with increasing flood risk. What should the assessment consider?}
	{\item[\textbf{(A)}] Nothing
		\item[\textbf{(B)}] Only the building's age
		\item[\textbf{(C)}] Only the color
		\item[\textbf{(D)}] Likelihood of flooding, resilience of facilities, insurance, and relocation or redundancy options}
	{D}
	{\item Physical events affect availability.
		\item Resilience lowers impact.
		\item Options guide investment.}
	{Include physical hazards in IT risk.}

	\hq{39}{Regulatory Fine Scenario}
	{How should potential regulatory penalties be included in a risk scenario?}
	{\item[\textbf{(A)}] As a component of impact, along with remediation costs and reputational damage
		\item[\textbf{(B)}] They should be ignored
		\item[\textbf{(C)}] They are the only impact
		\item[\textbf{(D)}] They are never estimable}
	{A}
	{\item Fines are one cost among many.
		\item Remediation can exceed fines.
		\item Reputation affects revenue.}
	{Include all impact types, not just fines.}

	\hq{40}{Open-Source Maintainer Risk}
	{A critical open-source component is maintained by one volunteer who has stopped responding. What is the risk?}
	{\item[\textbf{(A)}] No risk
		\item[\textbf{(B)}] Improved security
		\item[\textbf{(C)}] Faster updates
		\item[\textbf{(D)}] Unpatched vulnerabilities and lack of support, requiring contingency such as replacement or internal maintenance}
	{D}
	{\item Unsupported components accumulate risk.
		\item Inventory identifies dependence.
		\item Contingency plans reduce impact.}
	{Depend on a component, depend on its upkeep.}

	\hq{41}{Fourth-Party Concentration}
	{Several of your vendors all rely on the same small subcontractor. What does this reveal?}
	{\item[\textbf{(A)}] Diversification
		\item[\textbf{(B)}] No risk
		\item[\textbf{(C)}] Hidden concentration risk, because one failure could affect many vendors at once
		\item[\textbf{(D)}] Improved resilience}
	{C}
	{\item Shared dependencies concentrate risk.
		\item Visibility requires asking about subcontractors.
		\item Contingency planning is needed.}
	{Look for shared dependencies beyond direct vendors.}

	\hq{42}{Risk Proximity}
	{What does risk proximity describe?}
	{\item[\textbf{(A)}] How big the risk is
		\item[\textbf{(B)}] Who owns the risk
		\item[\textbf{(C)}] How soon a risk is expected to occur
		\item[\textbf{(D)}] How much the control costs}
	{C}
	{\item Proximity aids prioritization.
		\item It complements impact and likelihood.
		\item Near-term risks need quicker response.}
	{Proximity is about timing.}

	\hq{43}{After a Tool Change}
	{A key security tool is replaced by a new product. What should the risk practitioner do regarding risk ratings?}
	{\item[\textbf{(A)}] Keep ratings unchanged
		\item[\textbf{(B)}] Lower all ratings
		\item[\textbf{(C)}] Remove the controls
		\item[\textbf{(D)}] Reassess control effectiveness and coverage during and after the transition}
	{D}
	{\item Transitions often create gaps.
		\item New tools may behave differently.
		\item Evidence should confirm coverage.}
	{Change in controls means reassessment.}

	\hq{44}{End-User Computing}
	{Critical financial reports depend on complex spreadsheets maintained by individuals. What is the main risk?}
	{\item[\textbf{(A)}] No risk
		\item[\textbf{(B)}] Errors, lack of version control, and absence of access and change controls
		\item[\textbf{(C)}] Excess security
		\item[\textbf{(D)}] Too many backups}
	{B}
	{\item Spreadsheets lack inherent controls.
		\item Key-person dependency increases risk.
		\item Governance and testing reduce error.}
	{Critical spreadsheets need controls too.}

	\hq{45}{Industrial Control Assessment}
	{Which consideration is \textbf{MOST} important when assessing risk in industrial control systems?}
	{\item[\textbf{(A)}] Only confidentiality
		\item[\textbf{(B)}] Only website design
		\item[\textbf{(C)}] Only license cost
		\item[\textbf{(D)}] Safety and availability impacts, with limited ability to patch and test}
	{D}
	{\item Physical consequences can be severe.
		\item Patching may affect safety.
		\item Compensating controls are common.}
	{For OT, safety and availability lead.}

	\hq{46}{Assessing Remote Workforce Risk}
	{A company shifts to a largely remote workforce. Which factors should the assessment consider \textbf{FIRST}?}
	{\item[\textbf{(A)}] Endpoint security, home network exposure, access controls, data handling, and monitoring capability
		\item[\textbf{(B)}] Only office plants
		\item[\textbf{(C)}] Only commute times
		\item[\textbf{(D)}] Only desk size}
	{A}
	{\item Work location changes the threat surface.
		\item Endpoints and access are critical.
		\item Monitoring must adapt.}
	{Remote work changes the attack surface.}

	% ============================================================
	\section{Risk Response and Reporting (32 percent of the exam)}
	% ============================================================

	\hq{47}{Incident Playbooks}
	{Responders improvise during phishing incidents with inconsistent results. What is the \textbf{BEST} improvement?}
	{\item[\textbf{(A)}] Rely on individual judgment only
		\item[\textbf{(B)}] Delete the process
		\item[\textbf{(C)}] Develop and exercise playbooks with defined steps, roles, and decision points
		\item[\textbf{(D)}] Wait for audit}
	{C}
	{\item Playbooks make response repeatable.
		\item Exercises validate them.
		\item Roles reduce delay.}
	{Playbooks turn improvisation into process.}

	\hq{48}{Control Library}
	{Teams define controls differently for the same risks. What is the benefit of a central control library?}
	{\item[\textbf{(A)}] Consistent control definitions that can be mapped to risks and requirements and reused
		\item[\textbf{(B)}] More duplicate controls
		\item[\textbf{(C)}] Less clarity
		\item[\textbf{(D)}] No ownership}
	{A}
	{\item Reuse saves effort.
		\item Mapping supports compliance.
		\item Ownership is clearer.}
	{A shared library reduces duplication and confusion.}

	\hq{49}{Rationalizing Controls}
	{Control testing consumes many resources on low-value controls. What is the \textbf{BEST} approach?}
	{\item[\textbf{(A)}] Test everything equally
		\item[\textbf{(B)}] Focus testing on key controls tied to significant risks and reduce effort on low-risk ones
		\item[\textbf{(C)}] Stop testing
		\item[\textbf{(D)}] Test only the easiest}
	{B}
	{\item Risk-based testing optimizes resources.
		\item Key controls give the most assurance.
		\item Document the rationale.}
	{Test what matters most.}

	\hq{50}{Risk Acceptance Expiry}
	{Why should risk acceptances have an expiry date?}
	{\item[\textbf{(A)}] To annoy owners
		\item[\textbf{(B)}] To avoid acceptance
		\item[\textbf{(C)}] Conditions change, so acceptance must be revisited and reaffirmed or replaced by a response
		\item[\textbf{(D)}] To reduce documentation}
	{C}
	{\item Acceptance depends on assumptions.
		\item Expiry forces review.
		\item Unreviewed acceptances become stale.}
	{Acceptance should expire and be renewed consciously.}

	\hq{51}{Aging Issues}
	{A report shows many issues open far beyond their target dates. What is the \textbf{BEST} action?}
	{\item[\textbf{(A)}] Escalate overdue high-risk items, analyze root causes of delays, and reset realistic plans
		\item[\textbf{(B)}] Ignore the report
		\item[\textbf{(C)}] Delete old issues
		\item[\textbf{(D)}] Extend all dates silently}
	{A}
	{\item Aging means exposure persists.
		\item Root causes may be resourcing or priority.
		\item Realistic plans improve closure.}
	{Overdue issues need escalation and honest replanning.}

	\hq{52}{Vendor Remediation Escalation}
	{A vendor repeatedly misses remediation dates. Which action is \textbf{MOST} appropriate?}
	{\item[\textbf{(A)}] Escalate under contract terms, consider penalties or compensating controls, and reassess the relationship
		\item[\textbf{(B)}] Accept all delays
		\item[\textbf{(C)}] Remove the finding
		\item[\textbf{(D)}] Ignore it}
	{A}
	{\item Contracts provide remedies.
		\item Compensating controls reduce exposure.
		\item Reassessment may lead to exit.}
	{Use contractual levers for repeated delays.}

	\hq{53}{Service Credits}
	{A provider offers service credits for outages. How should risk be considered?}
	{\item[\textbf{(A)}] Credits eliminate the risk
		\item[\textbf{(B)}] Credits compensate cost but do not remove business impact, so continuity measures are still needed
		\item[\textbf{(C)}] Credits stop outages
		\item[\textbf{(D)}] Credits replace backups}
	{B}
	{\item Financial compensation is limited.
		\item Customers and operations still suffer.
		\item Resilience remains your responsibility.}
	{Credits are not resilience.}

	\hq{54}{KRIs for AI}
	{Which indicator is \textbf{MOST} useful for monitoring an AI system's risk?}
	{\item[\textbf{(A)}] Model accuracy trends, bias metrics, override rates, and complaint volumes
		\item[\textbf{(B)}] Number of servers only
		\item[\textbf{(C)}] Number of meetings
		\item[\textbf{(D)}] Marketing clicks}
	{A}
	{\item AI risks include drift and fairness.
		\item Human override shows trust and errors.
		\item Complaints reveal harm.}
	{Monitor accuracy, fairness, and overrides.}

	\hq{55}{Tuning DLP}
	{A DLP tool blocks many legitimate transfers and users bypass it. What is the \textbf{BEST} response?}
	{\item[\textbf{(A)}] Turn it off
		\item[\textbf{(B)}] Tune policies with business input, balance protection and productivity, and monitor exceptions
		\item[\textbf{(C)}] Block everything
		\item[\textbf{(D)}] Ignore user complaints}
	{B}
	{\item False positives drive circumvention.
		\item Business input improves rules.
		\item Exceptions need oversight.}
	{Controls that block the business get bypassed.}

	\hq{56}{Backup Test Results}
	{A restore test fails for a critical database. How should this be reported?}
	{\item[\textbf{(A)}] Hide it
		\item[\textbf{(B)}] Report only successes
		\item[\textbf{(C)}] As a significant control weakness with impact, root cause, remediation plan, and retest date
		\item[\textbf{(D)}] Consider it minor}
	{C}
	{\item Failure means data may be unrecoverable.
		\item Retest confirms the fix.
		\item Transparency supports decisions.}
	{Report failed restores clearly.}

	\hq{57}{Automating Control Evidence}
	{Collecting evidence for controls manually is slow and error-prone. What is the \textbf{BEST} improvement?}
	{\item[\textbf{(A)}] Automate evidence collection where feasible and validate its accuracy
		\item[\textbf{(B)}] Stop collecting evidence
		\item[\textbf{(C)}] Collect less often
		\item[\textbf{(D)}] Accept errors}
	{A}
	{\item Automation improves consistency.
		\item Validation ensures reliability.
		\item Frees time for analysis.}
	{Automate evidence to improve quality and speed.}

	\hq{58}{Board Cyber Report}
	{Which content is \textbf{MOST} useful in a board cyber risk report?}
	{\item[\textbf{(A)}] Raw logs
		\item[\textbf{(B)}] Tool brochures
		\item[\textbf{(C)}] Long technical glossaries
		\item[\textbf{(D)}] Top risks against appetite, trends, key incidents, control status, and decisions required}
	{D}
	{\item Boards need strategic information.
		\item Trends and appetite give context.
		\item Decisions drive action.}
	{Report risks, trends, and decisions.}

	\hq{59}{Incident Severity Levels}
	{Why define incident severity levels in advance?}
	{\item[\textbf{(A)}] To hide incidents
		\item[\textbf{(B)}] To avoid response
		\item[\textbf{(C)}] To delay decisions
		\item[\textbf{(D)}] To drive consistent response, escalation, and communication based on impact}
	{D}
	{\item Pre-defined levels speed decisions.
		\item Escalation paths align with severity.
		\item Consistency aids reporting.}
	{Define severity before the incident.}

	\hq{60}{Executive Communication During Incident}
	{During a major incident, how often should executives receive updates?}
	{\item[\textbf{(A)}] Only at the end
		\item[\textbf{(B)}] At agreed intervals defined by severity, even if there is no new information
		\item[\textbf{(C)}] Never
		\item[\textbf{(D)}] Whenever someone remembers}
	{B}
	{\item Regular updates reduce speculation.
		\item Cadence supports decisions.
		\item Silence erodes confidence.}
	{Predictable updates build trust.}

	\hq{61}{Cost of a Breach}
	{Which cost category is often \textbf{UNDERESTIMATED} in breach impact?}
	{\item[\textbf{(A)}] The price of coffee
		\item[\textbf{(B)}] Parking fees
		\item[\textbf{(C)}] Indirect costs such as lost customers, reputational damage, and long-term remediation
		\item[\textbf{(D)}] Stationery}
	{C}
	{\item Direct costs are easier to see.
		\item Indirect costs can be larger.
		\item Include them in scenarios.}
	{Do not forget the indirect costs.}

	\hq{62}{Recovery Priorities}
	{Many systems fail at once. How should recovery be prioritized?}
	{\item[\textbf{(A)}] Alphabetically
		\item[\textbf{(B)}] Randomly
		\item[\textbf{(C)}] By system age
		\item[\textbf{(D)}] By criticality to business processes, using BIA results and dependencies}
	{D}
	{\item BIA defines priorities.
		\item Dependencies set order.
		\item Business input resolves conflicts.}
	{Recover what the business needs first.}

	\hq{63}{Vendor Exit Execution}
	{A vendor relationship ends. Which activity is \textbf{MOST} important to manage risk?}
	{\item[\textbf{(A)}] Ignore the transition
		\item[\textbf{(B)}] Execute the exit plan: transfer data and services, revoke access, and verify deletion
		\item[\textbf{(C)}] Keep access open
		\item[\textbf{(D)}] Delete the contract only}
	{B}
	{\item Transition is a risky period.
		\item Access and data remain exposure.
		\item Verification provides evidence.}
	{Manage exit as a project, not an afterthought.}

	\hq{64}{Cyber Insurance Claim}
	{After an incident, the organization plans an insurance claim. What should it ensure?}
	{\item[\textbf{(A)}] Ignore the policy
		\item[\textbf{(B)}] Delay notification indefinitely
		\item[\textbf{(C)}] Evidence, timelines, and notification are handled according to policy terms
		\item[\textbf{(D)}] Destroy evidence}
	{C}
	{\item Policies have conditions.
		\item Late notice may affect coverage.
		\item Documentation supports the claim.}
	{Follow the policy's claim procedure.}

	\hq{65}{Learning From Near Misses}
	{Why record and analyze near misses?}
	{\item[\textbf{(A)}] They are meaningless
		\item[\textbf{(B)}] They create blame
		\item[\textbf{(C)}] They reveal control weaknesses and improvement opportunities without actual loss
		\item[\textbf{(D)}] They are illegal to record}
	{C}
	{\item Near misses are early warnings.
		\item Patterns suggest root causes.
		\item Learning is cheap compared with loss.}
	{Near misses are free lessons.}

	\hq{66}{Vulnerability Metrics Context}
	{A report states that vulnerabilities fell by 20 percent. What additional context is \textbf{MOST} needed?}
	{\item[\textbf{(A)}] The font used
		\item[\textbf{(B)}] Severity, exposure, asset criticality, and whether new scans covered the same scope
		\item[\textbf{(C)}] The author's name
		\item[\textbf{(D)}] The time of printing}
	{B}
	{\item Counts alone mislead.
		\item Scope changes affect comparisons.
		\item Critical assets matter more.}
	{Numbers need context.}

	\hq{67}{Patch Compliance Reporting}
	{A patch compliance report shows 98 percent, but the remaining 2 percent includes critical servers. What should the practitioner highlight?}
	{\item[\textbf{(A)}] That the exposure is concentrated in critical systems, which may represent significant risk
		\item[\textbf{(B)}] Only the overall 98 percent
		\item[\textbf{(C)}] Hide the exceptions
		\item[\textbf{(D)}] Consider the matter closed}
	{A}
	{\item Averages hide concentrated risk.
		\item Criticality matters.
		\item Highlight outliers.}
	{Do not let averages hide critical exceptions.}

	\hq{68}{Configuration Drift Reports}
	{What does recurring configuration drift in monitoring reports indicate?}
	{\item[\textbf{(A)}] Perfect compliance
		\item[\textbf{(B)}] No issues
		\item[\textbf{(C)}] Better security
		\item[\textbf{(D)}] Weak change control or automation, requiring root cause analysis and preventive measures}
	{D}
	{\item Drift means configurations are changing unmanaged.
		\item Preventive automation helps.
		\item Owners should be accountable.}
	{Recurring drift points to process weakness.}

	\hq{69}{Measuring Training Effectiveness}
	{Which indicator \textbf{BEST} demonstrates the effectiveness of security training?}
	{\item[\textbf{(A)}] Hours spent in class
		\item[\textbf{(B)}] Improved behavior metrics, such as reduced phishing click rates and increased reporting
		\item[\textbf{(C)}] Slides per course
		\item[\textbf{(D)}] Number of attendees}
	{B}
	{\item Behavior shows impact.
		\item Reporting rates show engagement.
		\item Attendance alone is weak.}
	{Measure what changes, not who attends.}

	\hq{70}{Phishing Simulation Use}
	{How should phishing simulation results be used?}
	{\item[\textbf{(A)}] To guide targeted education and measure trends, without punishing individuals
		\item[\textbf{(B)}] To fire staff who click
		\item[\textbf{(C)}] To publish names
		\item[\textbf{(D)}] To ignore results}
	{A}
	{\item Punishment reduces reporting.
		\item Targeted support improves behavior.
		\item Trends show progress.}
	{Use simulations to teach, not to shame.}

	\hq{71}{Response Budget Justification}
	{Which approach \textbf{BEST} justifies a budget request for risk response?}
	{\item[\textbf{(A)}] Present the risk addressed, expected reduction, cost, and alternatives in business terms
		\item[\textbf{(B)}] Fear alone
		\item[\textbf{(C)}] Technical jargon alone
		\item[\textbf{(D)}] No explanation}
	{A}
	{\item Decision makers need clear value.
		\item Alternatives show rigor.
		\item Business terms improve understanding.}
	{Make the case in business terms.}

	\hq{72}{Prioritizing Quick Wins}
	{Some risk responses are quick and cheap, others slow and expensive. How should they be sequenced?}
	{\item[\textbf{(A)}] Do only quick wins
		\item[\textbf{(B)}] Do only the slowest
		\item[\textbf{(C)}] Do nothing
		\item[\textbf{(D)}] Consider risk reduction, cost, and dependencies; do quick wins that matter while starting critical longer-term actions}
	{D}
	{\item Balance speed and value.
		\item Critical risks cannot wait.
		\item Dependencies affect ordering.}
	{Balance quick wins with critical work.}

	\hq{73}{Aligning Reports With Appetite}
	{Why present risk results relative to risk appetite?}
	{\item[\textbf{(A)}] It shows whether exposure is acceptable and where decisions are needed
		\item[\textbf{(B)}] It hides results
		\item[\textbf{(C)}] It replaces data
		\item[\textbf{(D)}] It avoids discussion}
	{A}
	{\item Appetite is the benchmark for acceptability.
		\item Gaps highlight needed action.
		\item It supports governance.}
	{Appetite gives risk numbers meaning.}

	\hq{74}{Risk Reporting Frequency}
	{How should reporting frequency be set?}
	{\item[\textbf{(A)}] The same for all risks always
		\item[\textbf{(B)}] Only annually
		\item[\textbf{(C)}] Based on risk volatility, audience needs, and defined escalation triggers
		\item[\textbf{(D)}] Only after incidents}
	{C}
	{\item Fast-moving risks need frequent updates.
		\item Audiences differ.
		\item Triggers handle exceptions.}
	{Frequency follows volatility and need.}

	\hq{75}{Verifying Remediation Claims}
	{A team claims that a vulnerability has been fixed. What should the practitioner do?}
	{\item[\textbf{(A)}] Accept without checking
		\item[\textbf{(B)}] Close immediately
		\item[\textbf{(C)}] Verify through rescanning or testing before closing the issue
		\item[\textbf{(D)}] Ignore it}
	{C}
	{\item Claims need evidence.
		\item Rescans confirm fixes.
		\item Verification prevents false assurance.}
	{Trust, then verify the fix.}

	\hq{76}{Communicating Residual Risk}
	{How should residual risk be communicated to decision makers?}
	{\item[\textbf{(A)}] Say it is zero
		\item[\textbf{(B)}] Hide it
		\item[\textbf{(C)}] Use only technical terms
		\item[\textbf{(D)}] Explain what remains after controls, its potential impact, and options to reduce it further}
	{D}
	{\item Residual risk is what they accept.
		\item Impact ties to objectives.
		\item Options support choice.}
	{Make residual risk visible and understandable.}

	\hq{77}{Metrics for Third-Party Risk Program}
	{Which metric \textbf{BEST} shows whether a third-party risk program is effective?}
	{\item[\textbf{(A)}] Percentage of critical vendors assessed on schedule, open high-risk findings, and time to remediate them
		\item[\textbf{(B)}] Number of vendors listed
		\item[\textbf{(C)}] Number of emails sent
		\item[\textbf{(D)}] Number of meetings held}
	{A}
	{\item Coverage of critical vendors matters.
		\item Open findings show exposure.
		\item Remediation time shows follow-through.}
	{Measure coverage, findings, and closure.}

	\hq{78}{Reporting Control Weaknesses Fairly}
	{How should control weaknesses be reported to management?}
	{\item[\textbf{(A)}] Exaggerate to get attention
		\item[\textbf{(B)}] Objectively, with evidence, impact, and recommended actions, without exaggeration or concealment
		\item[\textbf{(C)}] Hide minor ones always
		\item[\textbf{(D)}] Report opinions only}
	{B}
	{\item Credibility depends on accuracy.
		\item Evidence supports decisions.
		\item Recommendations enable action.}
	{Report objectively and completely.}

	% ============================================================
	\section{Information Technology and Security (22 percent of the exam)}
	% ============================================================

	\hq{79}{SASE}
	{Remote users and branch offices need consistent network security delivered from the cloud. Which model fits?}
	{\item[\textbf{(A)}] Secure access service edge, combining networking and security services in the cloud
		\item[\textbf{(B)}] A single office firewall only
		\item[\textbf{(C)}] Paper-based access
		\item[\textbf{(D)}] Dial-up modems}
	{A}
	{\item SASE converges network and security.
		\item Policies follow users.
		\item Cloud delivery scales.}
	{SASE delivers network and security as a cloud service.}

	\hq{80}{XDR}
	{What is the main advantage of extended detection and response?}
	{\item[\textbf{(A)}] Replacing all staff
		\item[\textbf{(B)}] Removing logs
		\item[\textbf{(C)}] Reducing visibility
		\item[\textbf{(D)}] Correlating telemetry across endpoints, network, email, and cloud for better detection and response}
	{D}
	{\item Cross-domain correlation reveals attack chains.
		\item Response can be coordinated.
		\item It needs tuning and skilled analysts.}
	{XDR connects signals across domains.}

	\hq{81}{CASB}
	{The organization needs visibility and policy control over the use of many cloud applications. Which technology helps?}
	{\item[\textbf{(A)}] A paper shredder
		\item[\textbf{(B)}] A cloud access security broker
		\item[\textbf{(C)}] A label printer
		\item[\textbf{(D)}] A battery backup}
	{B}
	{\item CASBs discover shadow cloud use.
		\item They enforce data and access policies.
		\item They integrate with identity.}
	{CASB governs cloud app usage.}

	\hq{82}{Hardware Security Module}
	{Why use a hardware security module?}
	{\item[\textbf{(A)}] To speed up printing
		\item[\textbf{(B)}] To store paper records
		\item[\textbf{(C)}] To cool servers
		\item[\textbf{(D)}] To generate, store, and use cryptographic keys in tamper-resistant hardware}
	{D}
	{\item Keys never leave the module in clear form.
		\item Tamper resistance protects against theft.
		\item Access is strongly controlled.}
	{HSMs protect keys in hardware.}

	\hq{83}{Infrastructure as Code}
	{Cloud infrastructure is defined in code templates. Which practice \textbf{BEST} reduces risk?}
	{\item[\textbf{(A)}] Allow unreviewed changes
		\item[\textbf{(B)}] Store templates on personal laptops
		\item[\textbf{(C)}] Scan templates for misconfigurations, use version control and peer review, and enforce policies in pipelines
		\item[\textbf{(D)}] Skip testing}
	{C}
	{\item Code errors can deploy insecure infrastructure at scale.
		\item Review and scanning catch issues early.
		\item Policy enforcement ensures consistency.}
	{Treat infrastructure code like application code.}

	\hq{84}{CI/CD Pipeline Security}
	{Which control \textbf{BEST} protects a CI/CD pipeline from compromise?}
	{\item[\textbf{(A)}] Least-privilege pipeline access, protected secrets, signed artifacts, and logging of changes
		\item[\textbf{(B)}] Shared pipeline passwords
		\item[\textbf{(C)}] Open access
		\item[\textbf{(D)}] No logs}
	{A}
	{\item Pipelines can push malicious code to production.
		\item Signing proves integrity.
		\item Logging supports investigation.}
	{Protect the build path to production.}

	\hq{85}{Kubernetes Risk}
	{A company runs workloads on Kubernetes. Which practice is \textbf{MOST} important?}
	{\item[\textbf{(A)}] Give all developers cluster admin
		\item[\textbf{(B)}] Run any image from the internet
		\item[\textbf{(C)}] Disable logging
		\item[\textbf{(D)}] Restrict cluster access, apply network policies, scan images, and secure configuration and secrets}
	{D}
	{\item Misconfiguration is a common issue.
		\item Network policies limit movement.
		\item Image and secret hygiene matter.}
	{Lock down the cluster.}

	\hq{86}{UEBA}
	{What does user and entity behavior analytics provide?}
	{\item[\textbf{(A)}] Faster printing
		\item[\textbf{(B)}] Physical access control
		\item[\textbf{(C)}] Detection of anomalous behavior by comparing activity to established baselines
		\item[\textbf{(D)}] Hardware cooling}
	{C}
	{\item Baselines reveal unusual activity.
		\item It helps detect compromised accounts and insiders.
		\item Tuning reduces false positives.}
	{Spot behavior that differs from normal.}

	\hq{87}{Deception Technology}
	{What is the main purpose of honeypots and deception technology?}
	{\item[\textbf{(A)}] To increase bandwidth
		\item[\textbf{(B)}] To replace firewalls
		\item[\textbf{(C)}] To lure attackers into decoys, detect their activity early, and learn their techniques
		\item[\textbf{(D)}] To store backups}
	{C}
	{\item Decoys generate high-confidence alerts.
		\item They reveal attacker behavior.
		\item They complement other controls.}
	{Decoys detect attackers early.}

	\hq{88}{BGP and Routing Risk}
	{Which risk is associated with weak internet routing security?}
	{\item[\textbf{(A)}] Faster internet always
		\item[\textbf{(B)}] Better privacy always
		\item[\textbf{(C)}] Fewer outages
		\item[\textbf{(D)}] Traffic could be misdirected or intercepted through route hijacking}
	{D}
	{\item Routing relies on trust between networks.
		\item Monitoring and validation reduce risk.
		\item Encryption protects content in transit.}
	{Routing can be hijacked; encrypt traffic.}

	\hq{89}{Email Gateway Security}
	{Which capability should a secure email gateway provide?}
	{\item[\textbf{(A)}] Blocking all email
		\item[\textbf{(B)}] Printing emails
		\item[\textbf{(C)}] Filtering for spam, malware, and phishing, plus attachment and link inspection
		\item[\textbf{(D)}] Deleting archives}
	{C}
	{\item Email is a primary attack vector.
		\item Inspection catches malicious content.
		\item Policy enforcement reduces risk.}
	{Filter the main attack path.}

	\hq{90}{Passkeys}
	{Which technology can reduce phishing risk by using cryptographic credentials tied to a device and website?}
	{\item[\textbf{(A)}] Longer typed passwords only
		\item[\textbf{(B)}] Security questions
		\item[\textbf{(C)}] Shared passwords
		\item[\textbf{(D)}] Passkeys using public-key cryptography}
	{D}
	{\item Credentials are bound to the legitimate site.
		\item No shared secret is sent.
		\item Devices protect the private key.}
	{Passkeys resist phishing.}

	\hq{91}{MFA Fatigue}
	{Attackers bombard a user with approval prompts until the user accepts one. Which control \textbf{BEST} counters this?}
	{\item[\textbf{(A)}] Allow unlimited prompts
		\item[\textbf{(B)}] Number matching or context in prompts, limits on prompts, and training
		\item[\textbf{(C)}] Disable MFA
		\item[\textbf{(D)}] Use SMS only}
	{B}
	{\item Prompt bombing exploits annoyance.
		\item Number matching makes approvals deliberate.
		\item Limits and monitoring detect abuse.}
	{Make MFA approvals deliberate.}

	\hq{92}{Password Managers}
	{Why might an organization encourage password managers?}
	{\item[\textbf{(A)}] They enable unique, strong passwords for each service without users memorizing them
		\item[\textbf{(B)}] They remove the need for any security
		\item[\textbf{(C)}] They reuse the same password
		\item[\textbf{(D)}] They disable MFA}
	{A}
	{\item Reuse is a major risk.
		\item Managers should be protected with MFA.
		\item They reduce weak password behavior.}
	{Unique passwords for every service.}

	\hq{93}{Prompt Injection}
	{A chatbot connected to company data follows malicious instructions hidden in user-provided content. What is this attack?}
	{\item[\textbf{(A)}] Disk failure
		\item[\textbf{(B)}] Cable damage
		\item[\textbf{(C)}] Screen flicker
		\item[\textbf{(D)}] Prompt injection}
	{D}
	{\item Models can follow untrusted instructions.
		\item Mitigation includes isolation and output controls.
		\item Limit the model's access and actions.}
	{Treat model inputs as untrusted.}

	\hq{94}{Model Theft and Leakage}
	{What is a risk of exposing a proprietary machine learning model through a public interface?}
	{\item[\textbf{(A)}] No risk
		\item[\textbf{(B)}] Better performance
		\item[\textbf{(C)}] Attackers could extract the model or sensitive training data through repeated queries
		\item[\textbf{(D)}] Lower cost always}
	{C}
	{\item Queries can reveal model behavior.
		\item Rate limits and monitoring help.
		\item Sensitive data should be minimized in training.}
	{Public models can leak what they learned.}

	\hq{95}{Log Management Architecture}
	{Which design \textbf{BEST} supports reliable log management?}
	{\item[\textbf{(A)}] Logs only on local devices
		\item[\textbf{(B)}] Logs editable by all
		\item[\textbf{(C)}] No retention
		\item[\textbf{(D)}] Centralized collection, time synchronization, protection from tampering, and defined retention}
	{D}
	{\item Central logs support correlation.
		\item Integrity supports investigations.
		\item Retention meets obligations.}
	{Centralize, protect, and retain logs.}

	\hq{96}{DNSSEC}
	{What does DNSSEC provide?}
	{\item[\textbf{(A)}] Authentication and integrity of DNS data, helping prevent spoofed responses
		\item[\textbf{(B)}] Encryption of all web content
		\item[\textbf{(C)}] Faster browsers
		\item[\textbf{(D)}] Spam filtering}
	{A}
	{\item Signatures validate DNS answers.
		\item It does not encrypt traffic.
		\item It reduces poisoning risk.}
	{DNSSEC authenticates DNS answers.}

	\hq{97}{PKI Trust}
	{Why is protecting a certificate authority's private key critical?}
	{\item[\textbf{(A)}] It only affects speed
		\item[\textbf{(B)}] It has no impact
		\item[\textbf{(C)}] Compromise would allow attackers to issue trusted fraudulent certificates
		\item[\textbf{(D)}] It affects printers}
	{C}
	{\item Trust chains depend on CA keys.
		\item Offline roots and HSMs are common protections.
		\item Revocation procedures must exist.}
	{CA keys are the root of trust.}

	\hq{98}{Data Backup Encryption}
	{Backups are stored offsite. Which control \textbf{BEST} protects them if lost or stolen?}
	{\item[\textbf{(A)}] Unencrypted storage
		\item[\textbf{(B)}] Encryption of backups with securely managed keys
		\item[\textbf{(C)}] Labels that say confidential
		\item[\textbf{(D)}] No access control}
	{B}
	{\item Encrypted backups are unreadable without keys.
		\item Key management is crucial.
		\item Access controls add protection.}
	{Encrypt the backups, protect the keys.}

	\hq{99}{Zero Trust Architecture Components}
	{Which combination \textbf{BEST} reflects a zero trust architecture?}
	{\item[\textbf{(A)}] Strong identity, device posture checks, least-privilege access, segmentation, and continuous monitoring
		\item[\textbf{(B)}] Network location as the only trust factor
		\item[\textbf{(C)}] Shared credentials
		\item[\textbf{(D)}] Perimeter only}
	{A}
	{\item Trust is evaluated per request.
		\item Context informs decisions.
		\item Monitoring detects misuse.}
	{Verify identity, device, and context continuously.}

	\hq{100}{Web Session Cookies}
	{Attackers steal session cookies to take over authenticated user sessions. Which controls \textbf{BEST} reduce this risk?}
	{\item[\textbf{(A)}] Long-lived cookies with no flags
		\item[\textbf{(B)}] Sending cookies over unencrypted connections
		\item[\textbf{(C)}] Sharing session IDs in URLs
		\item[\textbf{(D)}] Secure and HttpOnly cookie flags, TLS, short session lifetimes, and re-authentication for sensitive actions}
	{D}
	{\item Stolen sessions bypass passwords and MFA.
		\item Flags and TLS reduce theft opportunities.
		\item Short lifetimes limit misuse.}
	{Protect the session as carefully as the password.}

	% ============================================================
	\section{Conclusion}
	% ============================================================

	\begin{frame}{Exam Mindset Summary}
		\begin{block}{Remember These Core Principles}
			\begin{itemize}
				\item \textbf{Business First:} Risk is defined by its effect on business objectives.
				\item \textbf{Owners Decide:} Practitioners inform; owners accept, mitigate, transfer, or avoid.
				\item \textbf{Three Lines:} The first line owns risk, the second oversees, and the third gives independent assurance.
				\item \textbf{Ethics Matter:} Honesty, independence, and confidentiality underpin credible risk advice.
				\item \textbf{Evidence Over Opinion:} Use data, scenarios, and calibrated ratings.
				\item \textbf{Proportionality:} Controls should fit the risk and its cost.
				\item \textbf{Accountability Cannot be Outsourced:} Vendors and cloud providers share duties, but accountability remains with the organization.
			\end{itemize}
		\end{block}
	\end{frame}

	\begin{frame}{Thank You}
		\begin{center}
			{\Huge \textbf{Good Luck on Your CRISC Exam!}}

			\vspace{1cm}

			\normalsize
			Real exam questions test your ability to apply these concepts in practical, real-world scenarios.
		\end{center}
	\end{frame}

\end{document}
